\documentclass[11pt]{article}

\usepackage{amsmath, amssymb, amsthm}
\usepackage{bm}
\usepackage{geometry}
\usepackage{hyperref}
\usepackage{natbib}
\usepackage{xcolor}

\geometry{margin=1in}

\newcommand{\given}{\,|\,}
\newcommand{\Normal}{\mathcal{N}}
\newcommand{\Lognormal}{\text{Lognormal}}
\newcommand{\HalfNormal}{\text{HalfNormal}}
\newcommand{\Exponential}{\text{Exponential}}
\newcommand{\HalfCauchy}{\text{HalfCauchy}}
\newcommand{\Poisson}{\text{Poisson}}
\newcommand{\NegBin}{\text{NegBin}}

\setlength\parindent{0pt}

\title{\texttt{SCARCHhierarchSIR}\\Training and Forecasting Procedure}
\author{Tijs~W.~Alleman}
\date{\today}

\begin{document}
\maketitle


%-----------------------------------------------------------
\subsection*{Changes to \texttt{hierarchSIR}}
%-----------------------------------------------------------

The \texttt{SCARCHhierarchSIR} model extends the original \texttt{hierarchSIR} in two important ways. First, instead of treating U.S.\ states and territories as spatially independent units, we now incorporate spatial correlation in the Bayesian hierarchical inference layer. There is no spatial connection (e.g., mobility) between states through the forward-simulating SIR model. Forward simulations are performed using \texttt{diffrax}, which batches simulations across space and seasons using \texttt{jax} just-in-time compilation.\\

Second, the time-varying transmission modifier is generated by a spatially correlated AR(1)--GARCH(1,1) stochastic process whose hyperparameters are shared across seasons and states, whereas the \texttt{hierarchSIR} model's time-varying transmission modifiers were instantaneously mean-reverting (AR(1) with $\rho = 0$).


%-----------------------------------------------------------
\subsection*{Overview}
%-----------------------------------------------------------

The methodology consists of two distinct phases. \textbf{Training:} Embeds a forward-simulating SIR model in a Bayesian hierarchical inference framework. Season- and state-specific SIR parameters are modeled as exchangeable draws from shared hyperdistributions, whose posterior is inferred from historical influenza seasons. \textbf{Forecasting:} Uses the inferred hyperdistributions as informative priors for a Bayesian update of the within-season parameters during an ongoing season, producing four-week-ahead predictive distributions that are updated weekly as new data arrive.


%-----------------------------------------------------------
\subsection*{Training Dataset}
%-----------------------------------------------------------

A training dataset $\bm{D}$ spans at least one prior influenza season:
\begin{equation}
    \bm{D} = \bigl\{ D_{i,j}\ \forall\ i \in \bm{S},\ j \in \bm{T} \bigr\},
\end{equation}
Here, $\bm{S} = \{2023\text{-}2024,\, 2024\text{-}2025,\, 2025\text{-}2026\}$ is the set of training seasons, $\bm{T}$ is the set of U.S.\ states and territories, and $D_{i,j}$ is the NHSN HRD lab-confirmed hospital admissions for influenza in season $i$ and state $j$. We write $y_{i,j}(t)$ and $\hat{y}_{i,j}(t)$ for the observed and simulated counts respectively.

%-----------------------------------------------------------
\subsection*{Foward-simulating SIR model parameters}
%-----------------------------------------------------------

Three of the forward-simulating SIR parameters carry a non-centered hierarchy across seasons and U.S\ states and territories: the ascertainment rate $\rho$, the initial fraction of infected $f_I$ and recovered $f_R$. Opposed to \texttt{hierarchSIR}, the baseline transmission coefficient $\beta$ is fixed rather than inferred: $\beta_0 = 0.455$,  corresponding to a basic reproduction number of $R_0 = 1.6$.\\

The ascertainment rate in season $i$, U.S.\ state $j$ is modeled using the following log-additive structure,
\begin{align}
	 \log \rho_{i,j} &= \log \bar{\rho} + \sigma^{(i)}_\rho\, u^{(i)}_{\rho} + \sigma^{(j)}_\rho\, u^{(j)}_{\rho}, \nonumber\\
    \log \bar{\rho} &\sim \Normal\!\left(\mu_{\rho},1/9\right), \nonumber\\
    \sigma^{(s)}_\rho,\, \sigma^{(j)}_\rho &\sim \HalfNormal\! \left( 1/5 \right), \nonumber\\
    u^{(i)}_{\rho}, u^{(j)}_{\rho} &\sim \Normal(0,1),
\end{align}
where $\mu^{\rho}$ is obtained through a pre-optimization step. The initial fraction of infected $f_I$ is modeled identically. Because the values of $f_I$ are small ($\approx 10^{-5}$), a log-link structure was found to be computationally friendlier than the more natural logit-link structure.\\

The initial fraction of recovered $f_R$ in season $i$, U.S.\ state $j$ is modeled using a logit-additive structure to enforce the unit-interval constraint,
\begin{align}
	\text{logit}\,f_{R,i,j} &= \text{logit}\,\bar{f}_R + \sigma^{(i)}_{f_R}\, u^{(i)}_{f_R} + \sigma^{(j)}_{f_R}\, u^{(j)}_{f_R}, \nonumber\\
    \text{logit}\,\bar{f}_R &\sim \Normal\!\left(\text{logit}(0.4),\ 1\right), \nonumber\\
    \sigma^{(i)}_{f_R},\, \sigma^{(j)}_{f_R} &\sim \HalfNormal\!\left(1/5\right), \nonumber\\
    u^{(i)}_{f_R}, u^{(j)}_{f_R} &\sim \Normal(0,1),
\end{align}
The prior on $\text{logit}\,\bar{f}_R$ places the population-mean immunity near $40\pm10\%$, consistent with \texttt{hierarchSIR}.

%-----------------------------------------------------------
\subsection*{Spatially Correlated AR(1)--GARCH(1,1) Transmission Modifiers}\label{sec:argarch}
%-----------------------------------------------------------

The effective transmission coefficient in week $t$ in season $i$ and state $j$ in the SIR model is,
\begin{equation}
\beta_{i,j}(t) = \beta_0 \left[ 1 + \Delta \beta_{i,j}(t) \right],
\end{equation}
where the modifier trajectory $\Delta \beta_{i,j}(t)$ is modeled as a structured time series,
\begin{equation}
    \Delta \beta_{i,j} (t) = \Delta\bar{\beta_{j}}(t) + z_{i,j}(t),
\end{equation}
where $\Delta\bar{\beta_{j}}(t)$ is the seasonal mean modifier component, and $z_{i,j}(t)$ is a stochastic AR(1)--GARCH(1,1) component.

\paragraph{Seasonal mean modifiers}
A season-invariant, state-level, spatially correlated mean modifier is inferred for each week $t$ from mid October until mid April,
\begin{equation}
    \Delta\bar{\beta_{j}}(t) = 0.25\,\left[\bm{L}(\psi_1)^{-\top} \bm{v}(t)\right]_j, \qquad v_j(t) \sim \Normal(0,1),
\end{equation}
where the factor 0.25 sets the marginal prior standard deviation, and $\psi_1$ governs the spatial correlation strength with prior $\psi_1 \sim \text{Beta} \left( 3, 3\right)$.\\

Spatial correlation across U.S\ states and territories are encoded via a conditional autoregressive (CAR) precision matrix. Let $\bm{W}$ be the binary state-adjacency matrix with $W_{jj'} = 1$ if states $j$ and $j'$ share a border, and let $\bm{D}$ be the diagonal degree matrix. The precision matrix is computed as follows,
\begin{equation}
    \bm{Q}(\psi) = \bm{D} - \psi \bm{W}.
\end{equation}
The Cholesky factor $\bm{L}(\psi)$ of the precision matrix satisfies $\bm{Q}(\psi) = \bm{L}(\psi)\bm{L}(\psi)^\top$, and spatially correlated draws are obtained via $\bm{L}(\psi)^{-\top} \bm{u}$ with $\bm{u} \sim \Normal(\bm{0}, \bm{I})$ (Besag, 1974).

\paragraph{Stochastic component}
The stochastic part of $\Delta \beta_{i,j}(t)$, named $z_{i,j}(t)$, evolves according to an AR(1) process whose innovation variance follows a GARCH(1,1) model. At each week $t$ within a season, the recursion is,
\begin{alignat}{2}
	z_{i,j}(t) &= \phi_j \,z_{i,j}(t-1) + \varepsilon_{i,j}(t), \label{eq:ar1}\\
	\varepsilon_{i,j}(t) &= \eta_{i,j}(t)\,\sqrt{\sigma^2_{i,j}(t)}, \label{eq:garch_innov}\\
    \sigma^2_{i,j}(t) &= \omega + \alpha_j\,\varepsilon^2_{i,j}(t-1) + \beta_j\,\sigma^2_{i,j} (t-1), \label{eq:garch}
\end{alignat}
with initialization $z_{i,j}(0) = 0$ and $\varepsilon_{i,j}(0) = 0$. The driving noise $\bm{\eta}_j(t) = \{\eta_{i,j}(t)\}$ is spatially correlated across states:
\begin{equation}
    \bm{\eta}_j(t) = \bm{L}(\psi_{2})^{-\top}\, \bm{u}_j(t), \qquad u_{i,j}(t) \sim \Normal(0,1).
\end{equation}

The AR(1) persistence $\phi$ varies by state via a hierarchical logit-normal model,
\begin{align}
	\text{logit}~ \phi_j &= \text{logit}\,\bar{\phi} + \sigma^{(j)}_\phi\, u^{(j)}_{\phi}, \nonumber \\	
	\text{logit}\,\bar{\phi} &\sim \Normal(-1,\, 1), \nonumber\\
	\sigma^{(j)}_\phi &\sim \HalfNormal\!\left(1/2\right), \nonumber\\
	u^{(j)}_{\phi} &\sim \Normal(0,1).
\end{align}

The GARCH(1,1) steady-state baseline $\omega$ and total volatility persistence $\kappa_j = \alpha_j + \beta_j$ are modeled as,
\begin{align}
    \omega &\sim \Lognormal\!\left(\log(0.01/3),\ 1/25\right), \nonumber\\
    \text{logit}~\kappa_j &= \text{logit}\,\bar{\kappa} + \sigma^{(j)}_\kappa\, u^{(j)}_{\kappa},\nonumber\\
    \text{logit}\,\bar{\kappa} &\sim \Normal(-1,\, 1), \nonumber\\
    \sigma^{(j)}_\kappa &\sim \HalfNormal\!\left(1/2\right), \nonumber \\
    u^{(j)}_{\kappa} &\sim \Normal(0,1).
\end{align}
The total persistence of volatility is then split between the ARCH term $\alpha_j = \kappa_j \nu$ (sensitivity to new shocks) and the GARCH term $\beta_j = \kappa_j(1-\nu)$ (retention of volatility) via,
\begin{equation}
    \nu \sim \text{Beta}(5,\, 1),
\end{equation}
placing prior mass toward $\nu \approx 1$, i.e., incorporating the belief that when a modifier trajectory deviates substantially from the seasonal mean, the resulting increase in volatility is not retained long (short-term shock retention). The initial variance $\sigma^2_{i,j}(0)$ is centered at the GARCH(1,1) equilibrium variance of $\omega/(1-\kappa_j)$:
\begin{equation}
    \sigma^2_{i,j}(0) \sim \Lognormal\!\left(\log\!\left(\frac{\omega}{1-\kappa_s}\right),\ \sigma^2_{\sigma^2_0}\right), \qquad \sigma_{\sigma^2_0} \sim \HalfNormal\!\left(1/5\right).
\end{equation}

%-----------------------------------------------------------
\subsection*{Likelihood}
%-----------------------------------------------------------
 
The forward simulation model produces weekly predicted counts $\hat{y}_{i,j}(t)$, scaled to a 7-day observation window and passed through a softplus transformation, $\tilde{y} = \log(1 + e^y)$, to prevent negative values while remaining approximately linear for large positive outputs. A weighted negative binomial likelihood is used to accommodate overdispersion in hospital count data,
\begin{equation}
    \log \mathcal{L}(\bm{D} \given \bm{\Theta}) = \sum_{i}\sum_{j}\sum_{t}\, w_{i,j} \cdot \log p_{\NegBin}\!\left(y_{i,j}(t) \given \mu = \hat{y}_{i,j}(t),\, \alpha^{-1}_{j}\right),
\end{equation}
where $p_{\NegBin}(y \given \mu, \alpha^{-1})$ is the negative binomial probability mass function with mean $\mu$ and overdispersion parameter $\alpha^{-1}$. The overdispersion varies by state,
\begin{equation}
    \alpha^{-1}_j \sim \Lognormal\!\left(\log(0.0025),\ 1/25\right).
\end{equation}
The season-state weights $w_{i,j}$ normalize contributions so that no single large-epidemic season or high-burden state dominates the likelihood,
\begin{equation}
    w_{i,j} = \frac{\left(\max_{\{t\}} y_{i,j}(t)\right)^{-1}}{\dfrac{1}{|\bm{S}| \cdot |\bm{T}|} \displaystyle\sum_{i'}\sum_{j'} \left(\max_{\{t\}} y_{i',j'}(t)\right)^{-1}}.
\end{equation}
 
%-----------------------------------------------------------
\subsection*{Posterior and Inference}
%-----------------------------------------------------------
 
Collecting all within-season, within-state parameters into $\bm{\Theta}$ and all hyperparameters into $\bm{\Psi}$, the joint posterior during training is
\begin{equation}
    \mathcal{P}(\bm{\Theta}, \bm{\Psi} \given \bm{D}) \propto \mathcal{L}(\bm{D} \given \bm{\Theta})\, \mathcal{P}(\bm{\Theta} \given \bm{\Psi})\, \mathcal{P}(\bm{\Psi}),
\end{equation}
where $\mathcal{P}(\bm{\Theta} \given \bm{\Psi})$ encodes the hierarchical structure described above and $\mathcal{P}(\bm{\Psi})$ are the weakly informative hyperpriors listed in the preceding sections. Posterior sampling is performed using the No-U-Turn Sampler (NUTS) via PyMC.
 
%-----------------------------------------------------------
\subsection*{Forecasting}
%-----------------------------------------------------------
 
During an ongoing season, a reduced model is run in which all global and state-level components, as well as the season-level hyperparameters are replaced by the posterior median found during the training phase. Hence, only the current season's offsets $(u^{(i)}_\rho, u^{(i)}_{f_I}, u^{(i)}_{f_R})$ and the AR--GARCH stochastic trajectory $z_{i,j}(t)$ are sampled, conditioned on observations accumulated up to the forecast date.\\

%-----------------------------------------------------------
\bibliographystyle{plainnat}
\bibliography{references}
%-----------------------------------------------------------

\end{document}